
\baitracnghiem{thttlatex:01}{%601
Trong môi trường tabular, thông số nào sau đây cố định độ rộng một cột
}{
\bonpa{2}
{@$\{$2cm$\}$}
{p$\{$2cm$\}$}
{l}
{r}
}{ }

\baitracnghiem{thttlatex:02}{%602
Lệnh ngắt dòng nào sau đây khi ngắt dòng dồn đều các chữ còn lại trên dòng vừa ngắt:
}{%[2]
\bonpa{4}
{$\setminus$newline}
{$\backslash\backslash$ }
{$\backslash$cr}
{$\backslash$break}
}{ }

\baitracnghiem{thttlatex:03}{%603
Tệp ảnh có đuôi nào sau đây không dịch ra tệp PDF được khi đưa vào bằng lệnh\\
$\backslash$includegraphics[scale=1]$\{....\}$
}{%[2]
\bonpa{3}
{*.pdf}
{*.jpg}
{*.bmp}
{*.tif}
}{ }

\baitracnghiem{thttlatex:04}{%604
Lệnh nào sau đây luôn luôn định nghĩa một lệnh mới mà không thông báo lỗi:
}{%[2]
\bonpa{1}
{$\backslash$renewcommand}
{$\backslash$def}
{$\backslash$newcommand}
{$\backslash$newenvironment}
}{ }

\baitracnghiem{thttlatex:05}{%605
Lệnh nào sau đây đưa vào tiêu đề và mục lục nội dung khác nhau:
}{
\bonpa{2}
{$\backslash$section*[short-title]$\{$title$\}$}
{$\backslash$section[short-title]$\{$title$\}$}
{$\backslash$section$\{$title$\}$}
{$\backslash$section*$\{$title$\}$}
}{ }


\baitracnghiem{thttlatex:06}{%605
Lệnh nào sau đây chỉ đánh số trang và không tiêu đề chạy:
}{
\bonpa{3}
{$\backslash$pagestyle$\{$empty$\}$}
{$\backslash$pagestyle$\{$myheading$\}$}
{$\backslash$pagestyle$\{$plain$\}$}
{$\backslash$thispagestyle$\{$empty$\}$}
}{ }

\baitracnghiem{thttlatex:07}{%605
Nguyên tắc dấu trắng nào coi như không có dấu trắng:
}{
\bonpa{4}
{Hai hoặc nhiều dấu trắng liền nhau}
{Dấu cách (tabs)}
{Dấu xuống dòng}
{Các dấu trắng bắt đầu một dòng}
}{ }

\baitracnghiem{thttlatex:08}{%605
Muốn chủ động ngắt một từ ra làm đôi ở cuối dòng, thì ta làm cách nào sau đây: Ví dụ ngắt từ database thành data và base
}{
\bonpa{1}
{data$\backslash$-base}
{data$\backslash$newline base}
{data\~{}base}
{data base}
}{ }

\baitracnghiem{thttlatex:09}{%605
Một môi trường itemize như sau
$\backslash$begin$\{$itemize$\}$
$\backslash$item[(1)] Ví dụ
$\backslash$end$\{$itemize$\}$\\
khả năng nào là đúng
}{
\bonpa{2}
{$\bullet$ Ví dụ}
{(1) Ví dụ}
{$\bullet$ (1) Ví dụ}
{\textbf{Ví dụ}}
}{ }

\baitracnghiem{thttlatex:10}{%605
Để nội dung định lý không dùng chữ nghiêng, trước lệnh $\backslash $newtheorem$\{$theorem$\}\{$Định lý$\}$ người ta dùng lệnh nào sau đây:
}{
\bonpa{3}
{$\backslash$theoremstyle$\{$plain$\}$}
{$\backslash$theoremstyle$\{$main$\}$}
{$\backslash$theoremstyle$\{$definition$\}$}
{$\backslash$theoremstyle$\{$remark$\}$}
}{ }

\baitracnghiem{thttlatex:11}{%605
Những lệnh nào sau đây không dùng trong môi trường bảng:
}{
\bonpa{4}
{$\backslash$hline}
{$\backslash$cline$\{$1-3$\}$}
{$\backslash\backslash$[4pt]}
{$\backslash$vline}
}{ }

\baitracnghiem{thttlatex:12}{%605
Ký tự $\mathbb{R}$ là do lệnh nào dưới đây sinh ra:
}{
\bonpa{1}
{$\backslash$mathbb$\{$R$\}$}
{$\backslash$mathcal$\{$R$\}$}
{$\backslash$mathfrak$\{$R$\}$}
{$\backslash$mathscr$\{$R$\}$}
}{ }

\baitracnghiem{thttlatex:13}{%605
Đặt số trang dùng lệnh $\backslash$setcounter$\{<$số đếm$>\}\{$n$\}$, lệnh nào sau đây đặt lại số trang bằng 5:
}{
\bonpa{2}
{$\backslash$setcounter$\{$onpage$\}\{$5$\}$}
{$\backslash$setcounter$\{$page$\}\{$5$\}$}
{$\backslash$setcounter$\{\backslash $thepage$\}\{$5$\}$}
{$\backslash$setcounter$\{$lastpage$\}\{$5$\}$}
}{ }

\baitracnghiem{thttlatex:14}{%605
Mỗi dòng của trang văn bản được định nghĩa bằng 
$\backslash$setlength$\{\backslash$baselineskip$\}\{$14truept$\}$
Khi đó định nghĩa độ cao của trang bằng các lệnh sau, lệnh nào đặt không đúng.
}{
\bonpa{3}
{$\backslash$setlength$\{\backslash $textheight$\}\{$21truecm$\}$}
{$\backslash$setlength$\{\backslash $textheight$\}\{$20$\backslash$baselineskip$\}$}
{$\backslash$setlength$\{\backslash $textheight$\}\{$13$\}$}
{$\backslash$setlength$\{\backslash $textheight$\}\{$6in$\}$}
}{ }

\baitracnghiem{thttlatex:15}{%605
Lệnh gán nhãn
$\backslash$label$\{$nhanso$\}$
Lệnh nào sau đây dùng nhãn sai
}{
\bonpa{4}
{$\backslash$ref$\{$nhanso$\}$}
{$\backslash$pageref$\{$nhanso$\}$}
{$\backslash$eqref$\{$nhanso$\}$}
{$\backslash$cite$\{$nhanso$\}$}
}{ }

\baitracnghiem{thttlatex:16}{%605
Lệnh nào dùng thiết lập giao diện cho trình chiếu beamer sau đây:
}{
\bonpa{1}
{$\backslash$usetheme$\{$Antibes$\}$}
{$\backslash$usecolortheme$\{$Antibes$\}$}
{$\backslash$usefonttheme$\{$Antibes$\}$}
{$\backslash$usepackage$\{$Antibes$\}$}
}{ }

\baitracnghiem{thttlatex:17}{%605
Cặp lệnh nào sau đây làm ký hiệu biên tương ứng với nội dung bên trong công thức
}{
\bonpa{2}
{$\backslash$bigl( ...$\backslash$bigr)}
{$\backslash$left(...$\backslash$right)}
{$\backslash$Big(...$\backslash$Big)}
{$\backslash$Biggl(...$\backslash$Biggr)}
}{ }

\baitracnghiem{thttlatex:18}{%605
Để làm một ký tự đậm trong công thức toán, ta dùng lệnh nào sau đây:
}{
\bonpa{3}
{$\backslash$pmb$\{$...$\}$}
{$\backslash$bolsymbol$\{$...$\}$}
{$\backslash$mathbf$\{$...$\}$}
{$\backslash$textbf$\{$...$\}$}
}{ }

\baitracnghiem{thttlatex:19}{%605
Để cho một hình cố định tại vị trí đưa lệnh vào thì phần đầu tùy chọn nào hoặc không tùy chọn sau đây
}{
\bonpa{1}
{$\backslash$begin$\{$figure$\}[h]$}
{$\backslash$begin$\{$figure$\}[t]$}
{$\backslash$begin$\{$figure$\}[p]$}
{$\backslash$begin$\{$figure$\}$}
}{ }

\baitracnghiem{thttlatex:20}{%605
Đặt cỡ phông chữ cho văn bản $\backslash$documentcls[12pt]$\{$article $\}$, muốn tăng cỡ chữ lên 14pt ta thực hiện lệnh nào sau đây phù hợp nhất:
}{
\bonpa{1}
{Đặt lại $\backslash$documentcls[14pt]$\{$article$\}$}
{Dùng lệnh $\backslash$large}
{Dùng lệnh $\backslash$LARGE}
{Dùng lệnh $\backslash$huge}
}{ }

\baitracnghiem{thttlatex:21}{%605
Lệnh nào sau đây đánh lại số trang theo i,ii,iii,iv,....
}{
\bonpa{2}
{$\backslash$pagenumbering$\{$arabic$\}$}
{$\backslash$pagenumbering$\{$roman$\}$}
{$\backslash$pagenumbering$\{$alph$\}$}
{$\backslash$pagenumbering$\{$Alph$\}$}
}{ }

\baitracnghiem{thttlatex:22}{%605
Lệnh $\{\backslash$ allowdisplaybreaks ... $\}$ dùng ngắt khối nào sau đây
}{
\bonpa{2}
{Một bảng dài hơn một trang}
{Một khối dóng công thức cuối trang}
{Bất cứ một khối công thức nào dài quá 1 trang}
{Một hình lớn quá một trang}
}{ }

\baitracnghiem{thttlatex:23}{%605
Lệnh nào sau đây tách hai dòng trong một bảng
}{
\bonpa{3}
{$\backslash$hspace$\{$0.2cm$\}$}
{$\backslash$vspace$\{$0.2cm$\}$}
{$\backslash\backslash$[0.2cm]}
{$\backslash$cline$\{$0.2cm$\}$}
}{ }

\baitracnghiem{thttlatex:24}{%605
Lệnh đặc biệt nào chỉ dùng trong định nghĩa lệnh
}{
\bonpa{1}
{\#}
{\%}
{\&}
{@}
}{ }

\baitracnghiem{thttlatex:25}{%605
Làm lệnh phủ định một ký hiệu như $\not\subset$ ta dùng lệnh nào sau đây
}{
\bonpa{1}
{$\backslash$not$\{$...$\}$}
{$\backslash$overline$\{$...$\}$}
{$\backslash$overset$\{$/$\}$ $\{$...$\}$}
{$\backslash$stackrel$\{$...$\}$ $\{$/$\}$}
}{ }

\baitracnghiem{thttlatex:26}{%605
Môi trường nào sau đây dóng công thức nhiều dòng và thẳng hàng
}{
\bonpa{2}
{equation}
{align}
{multline}
{displaymath}
}{ }

\baitracnghiem{thttlatex:27}{%605
Dùng Geogebra khi Export ra tệp hình, sau đó đưa tệp hình vào văn bản TeX. Cách nào sau đây làm không đúng
}{
\bonpa{4}
{Đưa ra hình chọn đuôi *.png}
{Đưa ra tệp TeX, dùng standalone ra hình tệp *.pdf}
{Đưa ra tệp TeX, rồi đưa vào bằng $\backslash$input$\{$*.tex$\}$}
{Đưa ra tệp TeX, bỏ đầu, cuối, đưa vào $\backslash$input$\{$*.tex$\}$}
}{ }

\baitracnghiem{thttlatex:28}{%605
Nhận định nào về môi trường minipage thực hiện việc dóng song song sau đây là sai?
}{
\bonpa{2}
{dóng hai hình}
{dóng hai bảng}
{dóng một bảng và một hình}
{không dóng song song được}
}{ }

\baitracnghiem{thttlatex:29}{%605
Lệnh nào sau đây làm chữ trong dấu ngoặc nhọn không có chân như: \textsf{không chân}
}{
\bonpa{3}
{$\backslash $textrm $\{$...$\}$}
{$\backslash$ texttt $\{$...$\}$}
{$\backslash$ textsf $\{$...$\}$}
{$\backslash$ textit $\{$...$\}$}
}{ }

\baitracnghiem{thttlatex:30}{%605
Lệnh phép đồng dư nào sau đây làm ra : $x\equiv v \pmod{\theta}$
}{
\bonpa{3}
{\$x$\backslash$equiv v $\backslash$mod$\{\backslash $theta$\}$\$}
{\$x$\backslash $equiv v $\backslash$bmod$\{\backslash $theta$\}$\$}
{\$x$\backslash $equiv v $\backslash$pmod$\{\backslash $theta$\}$\$}
{\$x$\backslash $equiv v $\backslash $pod$\{\backslash $theta$\}$\$}
}{ }

